%% toc: P3 39 39.5
\section{Proofs: games for existential second-order expressibility}
\label{sec:annex-p03-w03-p03-058-eso-games}

Conventions are those of \cref{def:p03-w03-p03-058-setting}. Two facts about
existential second-order sentences are used repeatedly.

\emph{(F1) Normal form.} Every ESO-sentence is equivalent to one with
\(k\ge1\) quantified relation variables and a matrix of quantifier rank at
most \(r\) for any prescribed \(r\ge\operatorname{qr}(\psi)\), \(r\ge1\):
add a quantifier \((\exists P)\) over a unary variable \(P\) that does not
occur in \(\psi\). Since \(P\) does not occur in \(\psi\), an expansion
satisfies \(\psi\) whatever \(P\) denotes, and relations exist on every
nonempty universe (the empty one, for instance). A sentence of rank at most
\(\operatorname{qr}(\psi)\) has rank at most \(r\). The added variable is
unary, so monadic sentences stay monadic.

\emph{(F2) Finite Boolean combinations.} If \(\Psi_1,\ldots,\Psi_m\)
(\(m\ge1\)) are ESO-sentences \((\exists\bar P^{\,l})\psi_l\) whose tuples of
relation variables \(\bar P^{\,l}\) are pairwise disjoint (rename them apart
if necessary), then
\[
 \textstyle\bigwedge_l\Psi_l\ \text{ is equivalent to }\
 (\exists\bar P^{\,1})\cdots(\exists\bar P^{\,m})\bigwedge_l\psi_l,
 \qquad
 \bigvee_l\Psi_l\ \text{ is equivalent to }\
 (\exists\bar P^{\,1})\cdots(\exists\bar P^{\,m})\bigvee_l\psi_l .
\]
Each \(\psi_l\) mentions only \(\bar P^{\,l}\), so its truth in an expansion
depends only on the relations assigned to \(\bar P^{\,l}\). For the
conjunction, witnesses for the separate conjuncts together witness the
combined matrix, and conversely. For the disjunction, a witness for one
\(\Psi_l\), completed by arbitrary relations for the other variables, makes
the combined matrix true; conversely, if an expansion satisfies
\(\bigvee_l\psi_l\), some \(\psi_l\) holds there and its own relations
witness \(\Psi_l\). The combined sentences use relation variables of the
same arities as the \(\Psi_l\), so monadic sentences give monadic ones.

\begin{annexitem}{annex:p03-w03-p03-058-ef-constants}{thm:p03-w03-p03-058-ef-constants}{The first-order game with constants}
\begin{proof}
Write \(\sigma=\tau\cup\{c_1,\ldots,c_s\}\) with \(\tau\) the set of
relation symbols. For a \(\sigma\)-structure \(\mathbf A\) let
\(\mathbf A^\circ\) be its \(\tau\)-reduct and
\(\bar c^{\mathbf A}=(c_1^{\mathbf A},\ldots,c_s^{\mathbf A})\). The
reducts are structures for a finite relational signature, as in
\cref{def:p03-w03-p03-038-positions}.

\emph{The game.} A run of the \(r\)-move game on \(\mathbf A\) and
\(\mathbf B\) in the sense of \cref{def:p03-w03-p03-058-setting} is a play
of the pointed game
\(G_r(\mathbf A^\circ,\bar c^{\mathbf A};\mathbf B^\circ,\bar
c^{\mathbf B})\) of \cref{def:p03-w03-p03-038-positions}, with the same
moves. The winning conditions agree: the final position of length \(s+r\)
is a partial isomorphism exactly when the selected pairs together with the
constant pairs respect equality and preserve and reflect every relation
symbol, which is the condition of \cref{def:p03-w03-p03-058-setting}. So
strategies and winning strategies coincide, and by
\cref{thm:p03-w03-p03-038-pointed-ef} Duplicator wins exactly when
\((\mathbf A^\circ,\bar c^{\mathbf A})\equiv_r(\mathbf B^\circ,\bar
c^{\mathbf B})\), that is, when the pointed reducts agree on all
\(\tau\)-formulas of quantifier rank at most \(r\) with free variables
among \(x_1,\ldots,x_s\).

\emph{Translation.} Let \(\varphi\) be a \(\sigma\)-sentence. Rename its
bound variables so that none of \(x_1,\ldots,x_s\) is bound; the result is
equivalent to \(\varphi\) and has the same quantifier rank. Let
\(\varphi^*\) be obtained from it by replacing each \(c_j\) by \(x_j\). By
induction on subformulas \(\theta\), for every assignment \(t\) of their
other free variables,
\(\mathbf A\entails\theta[t]\iff\mathbf A^\circ\entails\theta^*[t\cup
\{x_j\mapsto c_j^{\mathbf A}\}]\): in atomic formulas a term \(c_j\) and the
variable \(x_j\) take the same value \(c_j^{\mathbf A}\); the Boolean steps
are immediate; and a quantified variable is never one of
\(x_1,\ldots,x_s\), so quantifying it commutes with fixing those values.
Hence \(\mathbf A\entails\varphi\iff\mathbf A^\circ\entails\varphi^*
[\bar c^{\mathbf A}]\) and \(\operatorname{qr}(\varphi^*)=
\operatorname{qr}(\varphi)\). Conversely, for a \(\tau\)-formula \(\chi\)
with free variables among \(x_1,\ldots,x_s\), let \(\chi[\bar c/\bar x]\) be
the result of replacing every free occurrence of \(x_j\) by \(c_j\). It is a
\(\sigma\)-sentence of the same quantifier rank, and the same induction
gives \(\mathbf A\entails\chi[\bar c/\bar x]\iff\mathbf A^\circ\entails
\chi[\bar c^{\mathbf A}]\).

\emph{Item 1.} If \((\mathbf A^\circ,\bar c^{\mathbf A})\equiv_r
(\mathbf B^\circ,\bar c^{\mathbf B})\) and \(\varphi\) is a
\(\sigma\)-sentence of rank at most \(r\), then
\(\mathbf A\entails\varphi\iff\mathbf A^\circ\entails\varphi^*[\bar
c^{\mathbf A}]\iff\mathbf B^\circ\entails\varphi^*[\bar c^{\mathbf B}]
\iff\mathbf B\entails\varphi\), so \(\mathbf A\equiv_r\mathbf B\). If
\(\mathbf A\equiv_r\mathbf B\) and \(\chi\) is a \(\tau\)-formula of rank at
most \(r\) with free variables among \(x_1,\ldots,x_s\), then
\(\mathbf A^\circ\entails\chi[\bar c^{\mathbf A}]\iff\mathbf A\entails
\chi[\bar c/\bar x]\iff\mathbf B\entails\chi[\bar c/\bar x]\iff
\mathbf B^\circ\entails\chi[\bar c^{\mathbf B}]\). So
\(\mathbf A\equiv_r\mathbf B\) holds exactly when the pointed reducts are
\(\equiv_r\)-equivalent, which holds exactly when Duplicator wins.

\emph{Item 2.} Let \(t=\operatorname{rec}_r(\mathbf A^\circ,\bar
c^{\mathbf A})\in\mathrm{Rec}_{s,r}\) (\cref{def:p03-w03-p03-038-records}),
let \(\chi_t\) be its characteristic formula
(\cref{def:p03-w03-p03-038-characteristic}), and put
\(\varphi^{\mathbf A}_r=\chi_t[\bar c/\bar x]\). By
\cref{thm:p03-w03-p03-038-characteristic} it is a \(\sigma\)-sentence of
quantifier rank at most \(r\), and since \(\mathrm{Rec}_{s,r}\) is finite
(\cref{thm:p03-w03-p03-038-finite-records}) the sentences
\(\varphi^{\mathbf A}_r\) form a finite set. For any \(\mathbf B\),
\(\mathbf B\entails\varphi^{\mathbf A}_r\iff\mathbf B^\circ\entails
\chi_t[\bar c^{\mathbf B}]\iff\operatorname{rec}_r(\mathbf B^\circ,\bar
c^{\mathbf B})=t\) by \cref{thm:p03-w03-p03-038-characteristic}, which by
\cref{thm:p03-w03-p03-038-pointed-ef} holds exactly when
\((\mathbf A^\circ,\bar c^{\mathbf A})\equiv_r(\mathbf B^\circ,\bar
c^{\mathbf B})\), that is, by item 1, when \(\mathbf A\equiv_r\mathbf B\).
Each \(\equiv_r\)-class is therefore the class of models of one of finitely
many sentences, so there are finitely many classes. When \(s=0\) the
positions are empty and the argument is unchanged.
\end{proof}
\end{annexitem}

Two consequences are used below. For \(\sigma\)-structures \(\mathbf X\) and
\(\mathbf Y\), let \(\Phi\) be the finite set of all sentences
\(\varphi^{\mathbf Z}_r\). Then
\(\mathbf X\equiv_r\mathbf Y\) if and only if \(\mathbf X\) and \(\mathbf Y\)
satisfy the same members of \(\Phi\): one direction holds because members of
\(\Phi\) have rank at most \(r\); for the other, \(\mathbf X\entails
\varphi^{\mathbf X}_r\) (as \(\mathbf X\equiv_r\mathbf X\)), so
\(\mathbf Y\entails\varphi^{\mathbf X}_r\) and
\(\mathbf X\equiv_r\mathbf Y\). In particular
\(\mathbf Y\entails\varphi^{\mathbf X}_r\) if and only if Duplicator wins the
\(r\)-move game on \(\mathbf X\) and \(\mathbf Y\).

\begin{annexitem}{annex:p03-w03-p03-058-preservation}{thm:p03-w03-p03-058-preservation}{Duplicator wins preserve existential second-order sentences}
\begin{proof}
Since \(\mathbf A\entails\Psi\), there are relations \(S_1,\ldots,S_k\) on
\(A\) of arities \(s_1,\ldots,s_k\) with
\((\mathbf A,S_1,\ldots,S_k)\entails\psi\). Let Spoiler choose them in the
relation-selection stage. Duplicator's winning strategy supplies
\(S'_1,\ldots,S'_k\) on \(B\) such that she wins the \(r\)-move game on
\((\mathbf A,S_1,\ldots,S_k)\) and \((\mathbf B,S'_1,\ldots,S'_k)\). These
are structures for the finite vocabulary \(\sigma\cup\{P_1,\ldots,P_k\}\),
so by \cref{thm:p03-w03-p03-058-ef-constants}(1) they satisfy the same
sentences of quantifier rank at most \(r\). As \(\psi\) is one of them,
\((\mathbf B,S'_1,\ldots,S'_k)\entails\psi\), and therefore
\(\mathbf B\entails\Psi\).
\end{proof}
\end{annexitem}

\begin{annexitem}{annex:p03-w03-p03-058-eso-method}{thm:p03-w03-p03-058-eso-method}{Soundness and completeness of the game method for ESO}
\begin{proof}
\emph{Soundness (item 1, ``if'').} Assume the condition and suppose, for a
contradiction, that an ESO-sentence \(\Psi\) defines \(Q\) on
\(\mathcal C\). By (F1) we may take \(\Psi=(\exists P_1)\cdots(\exists
P_k)\psi\) with \(k\ge1\), \(P_i\) of arity \(s_i\), and
\(\operatorname{qr}(\psi)\le r\) for some \(r\ge1\). The condition gives
\(\mathbf A,\mathbf B\in\mathcal C\) with \(Q(\mathbf A)=1\),
\(Q(\mathbf B)=0\) on which Duplicator wins the
\((\langle s_1,\ldots,s_k\rangle,r)\)-game. Then
\(\mathbf A\entails\Psi\), so \(\mathbf B\entails\Psi\) by
\cref{thm:p03-w03-p03-058-preservation}, so \(Q(\mathbf B)=1\): a
contradiction.

\emph{Completeness (item 1, ``only if'').} We prove the contrapositive.
Suppose that for some \(s_1,\ldots,s_k,r\) there is no pair
\(\mathbf A,\mathbf B\in\mathcal C\) with \(Q(\mathbf A)=1\),
\(Q(\mathbf B)=0\) on which Duplicator wins the
\((\langle s_1,\ldots,s_k\rangle,r)\)-game. Put
\(\sigma'=\sigma\cup\{P_1,\ldots,P_k\}\), \(P_i\) of arity \(s_i\), and let
\(\Phi\) be the finite set of sentences \(\varphi^{\mathbf X}_r\) for
\(\sigma'\)-structures \(\mathbf X\)
(\cref{thm:p03-w03-p03-058-ef-constants}(2)). For \(\mathbf A\in\mathcal C\)
let
\[
 \Theta(\mathbf A)=\{\varphi\in\Phi:(\mathbf A,S_1,\ldots,S_k)\entails
 \varphi\text{ for some }S_1,\ldots,S_k\text{ on }A\text{ of arities }
 s_1,\ldots,s_k\}.
\]
It is nonempty, since \(A\) carries at least one choice of relations. Let
\(\mathcal T\) be the set of subsets \(T\subseteq\Phi\) with
\(T=\Theta(\mathbf A)\) for some \(\mathbf A\in\mathcal C\) with
\(Q(\mathbf A)=1\); it is a finite set, being a set of subsets of the
finite set \(\Phi\). Let
\[
 \Psi=\bigvee_{T\in\mathcal T}\ \bigwedge_{\varphi\in T}
 (\exists P_1)\cdots(\exists P_k)\varphi ,
\]
read as the false sentence \((\exists x)\neg x=x\) if \(\mathcal T\) is
empty. By (F2) \(\Psi\) is equivalent to an ESO-sentence, and to a monadic
one when all \(s_i=1\).

If \(\mathbf A\in\mathcal C\) and \(Q(\mathbf A)=1\), then
\(\Theta(\mathbf A)\in\mathcal T\) and every \(\varphi\in\Theta(\mathbf A)\)
holds in some expansion of \(\mathbf A\), so \(\mathbf A\entails\Psi\).
Conversely let \(\mathbf B\in\mathcal C\) satisfy \(\Psi\). Then for some
\(\mathbf A\in\mathcal C\) with \(Q(\mathbf A)=1\), every member of
\(\Theta(\mathbf A)\) holds in some expansion of \(\mathbf B\). We show that
Duplicator wins the \((\langle s_1,\ldots,s_k\rangle,r)\)-game on
\(\mathbf A\) and \(\mathbf B\). Given Spoiler's relations
\(S_1,\ldots,S_k\) on \(A\), let \(\mathbf X=(\mathbf A,S_1,\ldots,S_k)\).
Then \(\varphi^{\mathbf X}_r\in\Theta(\mathbf A)\), so some
\(S'_1,\ldots,S'_k\) on \(B\) give
\((\mathbf B,S'_1,\ldots,S'_k)\entails\varphi^{\mathbf X}_r\); by the
consequence recorded after the proof of \cref{thm:p03-w03-p03-058-ef-constants},
Duplicator wins the \(r\)-move game on \(\mathbf X\) and
\((\mathbf B,S'_1,\ldots,S'_k)\). Choosing these \(S'_i\) and then playing
that strategy wins the game. By the hypothesis, \(Q(\mathbf B)\ne0\), so
\(Q(\mathbf B)=1\). Hence \(\Psi\) defines \(Q\) on \(\mathcal C\), and
\(Q\) is ESO-definable.

\emph{Item 2.} The same two arguments with all \(s_i=1\). In the soundness
direction a monadic \(\Psi\) is normalized by (F1) to a monadic sentence
with \(k\ge1\), and the game with arities \(\langle1,\ldots,1\rangle\) is
used. In the completeness direction the sentence \(\Psi\) built from
\(\langle1,\ldots,1\rangle\) is monadic by (F2).
\end{proof}
\end{annexitem}

\begin{annexitem}{annex:p03-w03-p03-058-np-conp}{thm:p03-w03-p03-058-np-conp}{\(\mathsf{NP}\ne\mathsf{coNP}\) as a statement about games}
\begin{proof}
Let \(\bar Q=1-Q\), a Boolean query on \(\mathcal G\). Statement (2) says
that for all \(s_1,\ldots,s_k,r\) there are \(\mathbf G,\mathbf H\in
\mathcal G\) with \(\bar Q(\mathbf G)=1\), \(\bar Q(\mathbf H)=0\) on which
Duplicator wins the \((\langle s_1,\ldots,s_k\rangle,r)\)-game. By
\cref{thm:p03-w03-p03-058-eso-method}(1) this holds exactly when \(\bar Q\) is
not ESO-definable on \(\mathcal G\). The class \(\{\mathbf G\in\mathcal G:
\bar Q(\mathbf G)=1\}\) is closed under isomorphism and \(\{E\}\) contains a
relation symbol, so by Fagin's theorem (\cref{thm:p03-fagin}) \(\bar Q\) is
ESO-definable on \(\mathcal G\) exactly when \(E(\bar Q)\in\mathsf{NP}\). It
therefore remains to show that \(E(\bar Q)\in\mathsf{NP}\) if and only if
\(\mathsf{NP}=\mathsf{coNP}\). Throughout, languages over \(\{0,1,\#\}\)
are read as binary languages through the code of the statement, and
complements are taken in \(\{0,1\}^*\).

\emph{Closure facts.} From \cref{def:p02-complexity-standing}: (a)
\(\mathsf P\) is closed under complement, since a deterministic decider can
reverse its answer, and \(\mathsf P\subseteq\mathsf{NP}\), with the empty
certificate. (b) If \(L\in\mathsf{NP}\), witnessed by \(q\) and \(R\), and
\(K\in\mathsf P\), then \(L\cap K\) and \(L\cup K\) are in \(\mathsf{NP}\):
use the predicates \(R(x,y)\wedge x\in K\) and \(R(x,y)\vee x\in K\), which
are polynomial-time, with the same bound \(q\). (c) If \(f\) is a
polynomial-time reduction from \(L\) to \(K\in\mathsf{NP}\), then
\(L\in\mathsf{NP}\): if \(f\) runs in time \(p(n)\), then
\(|f(x)|\le p(|x|)\), and with \(q\) taken nondecreasing (replace it by the
polynomial with the absolute values of its coefficients),
\(x\in L\) iff some \(y\) with \(|y|\le q(p(|x|))\) satisfies
\(R(f(x),y)\), a polynomial-time predicate of \((x,y)\).

\emph{Well-formed words.} Let \(W\) be the set of words \(a\#b\) in which
\(a\) is the binary representation of some \(n\ge1\) and \(b\in\{0,1\}^*\)
has length \(n^2\). These are exactly the encodings of graphs on
\(\{1,\ldots,n\}\) under \cref{def:p02-w02-p02-036-fagin-encoding}, so
\(E(Q)\) and \(E(\bar Q)\) are disjoint with union \(W\). Membership in
\(W\) is decided in polynomial time: check the shape, compute \(n^2\), which
has at most twice as many binary digits as \(a\), and compare it with the
length of \(b\). Hence
\[
 E(\bar Q)=W\setminus E(Q)=W\cap\overline{E(Q)},\qquad
 \overline{E(Q)}=E(\bar Q)\cup\overline W .
\]

If \(\mathsf{NP}=\mathsf{coNP}\), then \(\overline{E(Q)}\in\mathsf{coNP}=
\mathsf{NP}\), and \(E(\bar Q)=W\cap\overline{E(Q)}\in\mathsf{NP}\) by (b).

Conversely let \(E(\bar Q)\in\mathsf{NP}\). Then
\(\overline{E(Q)}=E(\bar Q)\cup\overline W\in\mathsf{NP}\) by (a) and (b).
Let \(L\in\mathsf{coNP}\). Then \(\overline L\in\mathsf{NP}\), and since
\(E(Q)\) is \(\mathsf{NP}\)-complete there is a polynomial-time \(f\) with
\(x\in\overline L\iff f(x)\in E(Q)\), that is,
\(x\in L\iff f(x)\in\overline{E(Q)}\). By (c), \(L\in\mathsf{NP}\). Thus
\(\mathsf{coNP}\subseteq\mathsf{NP}\). If \(L\in\mathsf{NP}\), then
\(\overline L\in\mathsf{coNP}\subseteq\mathsf{NP}\), so
\(L\in\mathsf{coNP}\); thus \(\mathsf{NP}\subseteq\mathsf{coNP}\), and the
classes are equal.

Combining: (2) holds iff \(E(\bar Q)\notin\mathsf{NP}\) iff
\(\mathsf{NP}\ne\mathsf{coNP}\).
\end{proof}
\end{annexitem}

\begin{annexitem}{annex:p03-w03-p03-058-three-col-np}{thm:p03-w03-p03-058-three-col-np}{3-colorability is \(\mathsf{NP}\)-complete}
\begin{proof}
Words over \(\{0,1,\#\}\) are read as binary words through the code of the
statement, and \(W\) is the polynomial-time set of well-formed encodings of
graphs from the proof of \cref{thm:p03-w03-p03-058-np-conp}.

\emph{Membership.} Let \(R(x,y)\) hold when \(x\in W\), say
\(x=e(\mathbf G)\) with universe \(\{1,\ldots,n\}\), and \(y\) has length
\(2n\) and lists \(n\) blocks from \(\{00,01,10\}\), read as colors
\(c(1),\ldots,c(n)\), with \(c(u)\ne c(v)\) whenever the digit of \(x\) for
the pair \((u,v)\) is \(1\). This is checked by \(n^2\) comparisons, and
\(n^2\) is at most the length of \(x\), so \(R\) is polynomial-time. A word
\(x\) lies in \(E(\mathrm{3COL})\) exactly when some \(y\) of length at most
twice the length of \(x\) satisfies \(R(x,y)\). So
\(E(\mathrm{3COL})\in\mathsf{NP}\) (\cref{def:p02-complexity-standing}).

\emph{The gadget.} In a graph with symmetric irreflexive edge relation,
let \(\mathsf T,\mathsf F,\mathsf N\) be three pairwise adjacent vertices,
and let \(a,b\) be vertices adjacent to \(\mathsf N\) (possibly \(a=b\)).
The \emph{gadget} \(\Gamma(a,b)\) adds three new vertices \(s,t,o\) and the
edges \(\{a,s\}\), \(\{b,t\}\), \(\{s,t\}\), \(\{s,o\}\), \(\{t,o\}\),
\(\{o,\mathsf N\}\); \(o\) is its \emph{output}. Fix a proper coloring \(c\)
of \(\mathsf T,\mathsf F,\mathsf N\), which uses three different colors, and
call \(c(\mathsf T)\) \emph{true} and \(c(\mathsf F)\) \emph{false}. Every
vertex adjacent to \(\mathsf N\) receives true or false in a proper
coloring. We claim:
\begin{enumerate}
\item[(G1)] in every proper coloring, if \(a\) and \(b\) are false, then
      \(o\) is false;
\item[(G2)] for any colors of \(a,b\) in \(\{c(\mathsf T),c(\mathsf F)\}\),
      if \(o\) is given the color true when \(a\) or \(b\) is true and false
      otherwise, then \(s\) and \(t\) can be colored so that the six gadget
      edges join different colors.
\end{enumerate}
For (G1): \(s\) and \(t\) differ from false, being adjacent to \(a\) and
\(b\), so \(\{c(s),c(t)\}=\{c(\mathsf T),c(\mathsf N)\}\), as they are
adjacent; then \(o\), adjacent to both, is false. For (G2), with \(o\) colored
as prescribed: if \(a,b\) are false, put \(s\) true and \(t\) the color of
\(\mathsf N\); if \(a\) is true and \(b\) false, put \(s\) false and \(t\) the
color of \(\mathsf N\); if \(a\) is false and \(b\) true, put \(s\) the color
of \(\mathsf N\) and \(t\) false; if both are true, put \(s\) false and \(t\)
the color of \(\mathsf N\). In each case \(s\ne a\), \(t\ne b\), \(s\ne t\),
\(s\ne o\), \(t\ne o\) and \(o\ne c(\mathsf N)\), as a check of the four
cases shows (in the first case \(o\) is false, in the others true).

\emph{From clause sets to graphs.} Let \(S\) be a finite set of nonempty
clauses, each with at most three literals, over a finite set \(U\) of
variables. The graph \(H_S\) has the vertices \(\mathsf T,\mathsf F,
\mathsf N\), pairwise adjacent; for each \(w\in U\) two vertices \(w^+,w^-\)
forming a triangle with \(\mathsf N\); and for each clause, listed as
\(\{l_1,l_2,l_3\}\) with a literal repeated if the clause has fewer than three,
the gadget \(\Gamma(\bar l_1,\bar l_2)\) with output \(o_1\), the gadget
\(\Gamma(o_1,\bar l_3)\) with output \(o_2\), and the edge
\(\{o_2,\mathsf F\}\). Here \(\bar l=w^+\) for the literal \(w\) and
\(\bar l=w^-\) for \(\neg w\). All edges are taken in both directions, and
each gadget has its own new vertices.

\emph{\(S\) is satisfiable iff \(H_S\) is 3-colorable.} Let \(c\) be a proper
coloring. Define \(v(w)=\mathbf T\) when \(w^+\) is true. Since \(w^+\),
\(w^-\) and \(\mathsf N\) form a triangle, exactly one of \(w^+,w^-\) is true,
so \(\bar l\) is true exactly when \(v\) makes the literal \(l\) true. In each
clause gadget pair, \(o_2\) is adjacent to \(\mathsf F\) and \(\mathsf N\),
so it is true. By (G1) for \(\Gamma(o_1,\bar l_3)\), \(o_1\) or \(\bar l_3\)
is true; if \(o_1\) is true, by (G1) for \(\Gamma(\bar l_1,\bar l_2)\) one of
\(\bar l_1,\bar l_2\) is true. So \(v\) makes a literal of each clause true.
Conversely let \(v\) satisfy \(S\). Color \(\mathsf T,\mathsf F,\mathsf N\)
by \(1,2,3\); color \(w^+\) true and \(w^-\) false when \(v(w)=\mathbf T\),
and the other way round otherwise; color \(o_1\) true exactly when \(v\) makes
\(l_1\) or \(l_2\) true, and \(o_2\) true exactly when \(o_1\) or \(\bar l_3\)
is true, which holds because the clause is satisfied; and color the inner
vertices of each gadget by (G2). The triangles get three colors; every gadget
edge is proper by (G2); and \(\{o_2,\mathsf F\}\) is proper because \(o_2\) is
true. The inner vertices \(s,t\) of a gadget are adjacent only to vertices of
that gadget and to its inputs, whose colors were fixed before the inner
vertices were chosen, so the separate choices are compatible.

\emph{Hardness.} Let \(L\in\mathsf{NP}\). By \cref{thm:p02-sat-np-complete}
there is a polynomial-time reduction \(f\) from \(L\) to \(\mathsf{SAT}\)
(\cref{def:p02-sat-taut}). Define \(g\) on binary words: if \(z\) encodes a
formula \(\varphi\), let \(g(z)\) be the encoding \(e(H_{S_\varphi})\), with
the vertices numbered \(1,\ldots,n\) in any fixed order, where
\(S_\varphi\) is the clause set of \cref{def:p02-tseitin-construction};
otherwise let \(g(z)=e(K_4)\), the complete graph on four vertices, which is
not 3-colorable because four pairwise adjacent vertices need four colors. All
clauses of \(S_\varphi\) are nonempty with at most three literals
(\cref{thm:p02-tseitin-size}), and \(S_\varphi\) is satisfiable iff
\(\varphi\) is (\cref{thm:p02-tseitin}). So \(z\in\mathsf{SAT}\) iff
\(H_{S_\varphi}\) is 3-colorable iff \(g(z)\in E(\mathrm{3COL})\), and a
malformed \(z\) lies in neither language. As for computing \(g\): recognizing
formula encodings is linear-time (\cref{def:p02-sat-taut}); \(S_\varphi\) is
produced in polynomial time and has at most \(3M+1\) clauses over at most
\(2N\) variables (\cref{thm:p02-tseitin-size}); so \(H_{S_\varphi}\) has
\(n\le3+4N+6(3M+1)\) vertices, polynomial in the length of \(z\), and its
adjacency word of length \(n^2\) is written by listing the
\(O(N+M)\) edges and filling an \(n\times n\) table. Hence \(g\) is
polynomial-time, and \(g\circ f\) reduces \(L\) to \(E(\mathrm{3COL})\) in
polynomial time (\cref{thm:p02-w02-p02-019-reductions}).
\end{proof}
\end{annexitem}

\begin{annexitem}{annex:p03-w03-p03-058-np-conp-3col}{thm:p03-w03-p03-058-np-conp-3col}{Kolaitis's game form of \(\mathsf{NP}\ne\mathsf{coNP}\)}
\begin{proof}
By \cref{thm:p03-w03-p03-058-three-col-np}, \(Q=\mathrm{3COL}\) satisfies
the hypothesis of \cref{thm:p03-w03-p03-058-np-conp}, whose item (2) for this
\(Q\) is the game statement here.
\end{proof}
\end{annexitem}

\begin{annexitem}{annex:p03-w03-p03-058-af-method}{thm:p03-w03-p03-058-af-method}{Soundness and completeness of Ajtai--Fagin games}
\begin{proof}
\emph{Item 3.} Let \(\mathbf A,\mathbf B\) be as in the hypothesis.
Duplicator chooses \(\mathbf A\). Whatever relations \(S_1,\ldots,S_k\)
Spoiler chooses on \(A\), she chooses \(\mathbf B\) and the relations
\(S'_1,\ldots,S'_k\) that her winning strategy in the
\((\langle s_1,\ldots,s_k\rangle,r)\)-Ehrenfeucht--Fra\"iss\'e game on
\(\mathbf A\) and \(\mathbf B\) prescribes, and continues with that
strategy in the \(r\)-move game.

\emph{Soundness of items 1 and 2.} Suppose \(Q\) is ESO-definable on
\(\mathcal C\) (monadic ESO-definable, for item 1). By (F1) it is defined by
\(\Psi=(\exists P_1)\cdots(\exists P_k)\psi\) with \(k\ge1\), \(P_i\) of
arity \(s_i\) (all \(s_i=1\) in the monadic case) and
\(\operatorname{qr}(\psi)\le r\), \(r\ge1\). We show that Duplicator does
not win the \((\langle s_1,\ldots,s_k\rangle,r)\)-Ajtai--Fagin game for
\(Q\) on \(\mathcal C\). Let \(\mathbf A\in\mathcal C\) with
\(Q(\mathbf A)=1\) be any first move of Duplicator (if there is none, she
cannot move and loses). Then \(\mathbf A\entails\Psi\); let Spoiler choose
\(S_1,\ldots,S_k\) with \((\mathbf A,S_1,\ldots,S_k)\entails\psi\). If
Duplicator could choose \(\mathbf B\in\mathcal C\) with \(Q(\mathbf B)=0\)
and \(S'_1,\ldots,S'_k\) with a win in the \(r\)-move game on the expanded
structures, then by \cref{thm:p03-w03-p03-058-ef-constants}(1)
\((\mathbf B,S'_1,\ldots,S'_k)\entails\psi\), hence
\(\mathbf B\entails\Psi\) and \(Q(\mathbf B)=1\), a contradiction. So for
these parameters Duplicator does not win; contraposing, if she wins for all
parameters, \(Q\) is not (monadic) ESO-definable.

\emph{Completeness of items 1 and 2.} Suppose that for some
\(s_1,\ldots,s_k,r\) (all \(s_i=1\) for item 1) Duplicator does not win the
Ajtai--Fagin game. By the definition of winning, for every
\(\mathbf A\in\mathcal C\) with \(Q(\mathbf A)=1\) there are relations
\(S_1,\ldots,S_k\) on \(A\) such that for no \(\mathbf B\in\mathcal C\) with
\(Q(\mathbf B)=0\) and no \(S'_1,\ldots,S'_k\) on \(B\) does Duplicator win
the \(r\)-move game on \((\mathbf A,S_1,\ldots,S_k)\) and
\((\mathbf B,S'_1,\ldots,S'_k)\). With \(\sigma'\) and \(\Phi\) as in the
proof of \cref{thm:p03-w03-p03-058-eso-method}, let \(\Phi_0\) be the set of
\(\varphi\in\Phi\) such that
\begin{itemize}
\item some \(\mathbf A\in\mathcal C\) with \(Q(\mathbf A)=1\) has an
      expansion \((\mathbf A,S_1,\ldots,S_k)\entails\varphi\), and
\item no \(\mathbf B\in\mathcal C\) with \(Q(\mathbf B)=0\) has an expansion
      \((\mathbf B,S'_1,\ldots,S'_k)\entails\varphi\),
\end{itemize}
and let \(\Psi=(\exists P_1)\cdots(\exists P_k)\bigvee\Phi_0\), with the
false sentence \((\exists x)\neg x=x\) as matrix if \(\Phi_0\) is empty.
This is an ESO-sentence with relation variables of arities
\(s_1,\ldots,s_k\), monadic for item 1.

Let \(\mathbf A\in\mathcal C\) with \(Q(\mathbf A)=1\), and take
\(S_1,\ldots,S_k\) as above; put \(\mathbf X=(\mathbf A,S_1,\ldots,S_k)\)
and \(\varphi=\varphi^{\mathbf X}_r\). Then \(\mathbf X\entails\varphi\).
If some \(\mathbf B\in\mathcal C\) with \(Q(\mathbf B)=0\) had an expansion
\(\mathbf Y\entails\varphi\), Duplicator would win the \(r\)-move game on
\(\mathbf X\) and \(\mathbf Y\) (the consequence recorded after
the proof of \cref{thm:p03-w03-p03-058-ef-constants}), contrary to the choice of
\(S_1,\ldots,S_k\). So \(\varphi\in\Phi_0\) and \(\mathbf A\entails\Psi\).
Conversely, if \(\mathbf B\in\mathcal C\) satisfies \(\Psi\), some expansion
of \(\mathbf B\) satisfies some \(\varphi\in\Phi_0\), and by the second
clause in the definition of \(\Phi_0\), \(Q(\mathbf B)\ne0\). So \(\Psi\)
defines \(Q\) on \(\mathcal C\).
\end{proof}
\end{annexitem}

The next two proofs concern \(k\)-coloured graphs
(\cref{def:p03-w03-p03-058-neighbourhoods}). A \emph{local isomorphism}
from \(U\subseteq A\) to \(V\subseteq B\) is a bijection \(h\colon U\to V\)
that preserves colours and satisfies \(E(u,u')\iff E(h(u),h(u'))\) for
\(u,u'\in U\); it is an isomorphism between the induced substructures, and
its restriction to any subset of \(U\) is a local isomorphism onto the
image.

\emph{(L) Transport of neighbourhoods.} Let \(h\colon U\to V\) be a local
isomorphism, \(u\in U\) and \(t\ge1\), with \(N(u,t)\subseteq U\) and
\(N(h(u),t)\subseteq V\). Then \(h\) maps \(N(u,t)\) onto \(N(h(u),t)\) and
\(\operatorname{dist}(h(u),h(x))=\operatorname{dist}(u,x)\) for
\(x\in N(u,t)\).

\emph{Proof of (L).} By induction on \(t'\le t\) we show that \(h\) maps
the set of points at distance exactly \(t'-1\) from \(u\) onto the set of
points at distance exactly \(t'-1\) from \(h(u)\). For \(t'=1\) both sets
are singletons \(\{u\}\) and \(\{h(u)\}\). Let \(t'<t\) and \(x\) be at
distance \(t'\) from \(u\). On a shortest path, \(x\) has a neighbour
\(y\) at distance \(t'-1\); both lie in \(N(u,t)\subseteq U\), so
\(E(h(y),h(x))\), and \(h(y)\) is at distance \(t'-1\) from \(h(u)\) by
induction. Hence \(\operatorname{dist}(h(u),h(x))\le t'\). It is not
smaller: otherwise \(h(x)\) is at distance \(<t'\) from \(h(u)\), so by
induction \(h(x)=h(z)\) for some \(z\) at distance \(<t'\) from \(u\), and
injectivity gives \(x=z\). The hypotheses are symmetric in \(h\) and
\(h^{-1}\), so the same argument applied to \(h^{-1}\) shows that every point
at distance \(t'\) from \(h(u)\) is the image of a point at distance \(t'\)
from \(u\). \qquad\(\square\)

\emph{(T) Smaller radii.} If \(x\in A\) and \(y\in B\) have the same
\(d\)-type and \(1\le\rho\le d\), they have the same \(\rho\)-type: the
isomorphism \(g\) of the \(d\)-neighbourhoods with \(g(x)=y\) is a local
isomorphism from \(N(x,d)\) onto \(N(y,d)\), and (L) with \(t=\rho\) shows
that it maps \(N(x,\rho)\) onto \(N(y,\rho)\), the restriction being an
isomorphism of the \(\rho\)-neighbourhoods sending \(x\) to \(y\).
Consequently the \(\rho\)-type of an element is a function of its
\(d\)-type, and the number of elements of a given \(\rho\)-type is the sum of
the numbers of elements of the \(d\)-types that determine it. So finite
\(d\)-equivalent coloured graphs are \(\rho\)-equivalent for every
\(\rho\le d\).

\begin{annexitem}{annex:p03-w03-p03-058-hanf}{thm:p03-w03-p03-058-hanf}{Counting neighbourhood types wins the first-order game}
\begin{proof}
Put \(\rho_j=3^{r-j}\) for \(0\le j\le r\). After \(j\) moves, with
\(\bar a=(a_1,\ldots,a_j)\) selected in \(\mathbf A\) and
\(\bar b=(b_1,\ldots,b_j)\) in \(\mathbf B\), Duplicator maintains:

\((I_j)\) there is a local isomorphism \(h\) from \(N(\bar a,\rho_j)\) onto
\(N(\bar b,\rho_j)\) with \(h(a_i)=b_i\) for \(i\le j\).

\((I_0)\) holds with the empty map. If \((I_r)\) holds, then since
\(\rho_r=1\) and \(N(a_i,1)=\{a_i\}\), \(h\) is a local isomorphism from
\(\{a_1,\ldots,a_r\}\) onto \(\{b_1,\ldots,b_r\}\) sending \(a_i\) to
\(b_i\). Its graph \(\{(a_i,b_i)\}\) is therefore a partial isomorphism: it
is a function and injective because \(h\) is, and it preserves and
reflects \(E\) and every colour. So it suffices to show that \((I_j)\) can
be extended to \((I_{j+1})\) whatever Spoiler plays, for \(j<r\).

Assume \((I_j)\) with map \(h\), and put \(\rho=\rho_j\),
\(\rho'=\rho_{j+1}\), so \(\rho=3\rho'\). For each \(i\le j\),
\(N(a_i,\rho)\) and \(N(b_i,\rho)\) lie in the domain and range of \(h\), so
by (L) \(h\) maps \(N(a_i,t)\) onto \(N(b_i,t)\) for every \(t\le\rho\) and
preserves distances from \(a_i\) there. Suppose Spoiler plays \(a\in A\); a
move in \(\mathbf B\) is handled by the same argument with the roles of the
structures exchanged and \(h^{-1}\) in place of \(h\), since all hypotheses
are symmetric.

\emph{Case 1: \(\operatorname{dist}(a,a_i)\le2\rho'-1\) for some \(i\).}
Let \(b=h(a)\); this is defined because \(2\rho'-1<\rho\). Then
\(\operatorname{dist}(b_i,b)=\operatorname{dist}(a_i,a)\). Every \(x\) with
\(\operatorname{dist}(a,x)\le\rho'-1\) has
\(\operatorname{dist}(a_i,x)\le3\rho'-2<\rho\), so
\(N(a,\rho')\subseteq N(a_i,\rho)\), and likewise
\(N(b,\rho')\subseteq N(b_i,\rho)\). By (L), \(h\) maps \(N(a,\rho')\) onto
\(N(b,\rho')\), and it maps each \(N(a_l,\rho')\) onto \(N(b_l,\rho')\).
The restriction of \(h\) to \(N(\bar aa,\rho')\) is therefore a local
isomorphism onto \(N(\bar bb,\rho')\) sending \(a_l\) to \(b_l\) and \(a\)
to \(b\), which is \((I_{j+1})\).

\emph{Case 2: \(\operatorname{dist}(a,a_i)\ge2\rho'\) for every \(i\)}
(this includes \(j=0\)). Let \(\theta\) be the \(\rho'\)-type of \(a\), and
let \(X_{\mathbf A}\) be the set of elements of \(A\) of \(\rho'\)-type
\(\theta\) at distance at most \(2\rho'-1\) from some \(a_i\), and
\(X_{\mathbf B}\) the corresponding set in \(B\) with respect to \(\bar b\).
If \(x\in X_{\mathbf A}\), with \(\operatorname{dist}(x,a_i)\le2\rho'-1\),
then exactly as in Case 1 \(h\) maps \(N(x,\rho')\) isomorphically onto
\(N(h(x),\rho')\) with \(x\mapsto h(x)\), so \(h(x)\) has \(\rho'\)-type
\(\theta\), and \(\operatorname{dist}(h(x),b_i)=\operatorname{dist}(x,a_i)
\le2\rho'-1\); thus \(h(x)\in X_{\mathbf B}\). The same argument for
\(h^{-1}\) shows \(h^{-1}(X_{\mathbf B})\subseteq X_{\mathbf A}\), so
\(|X_{\mathbf A}|=|X_{\mathbf B}|\). Since \(\rho'\le\rho_1=3^{r-1}\le d\),
the structures are \(\rho'\)-equivalent by (T), and both are finite, so the
numbers of elements of type \(\theta\) outside \(X_{\mathbf A}\) and outside
\(X_{\mathbf B}\) are equal. In \(\mathbf A\) this number is at least
\(1\), because \(a\notin X_{\mathbf A}\). Choose \(b\) of type \(\theta\)
outside \(X_{\mathbf B}\); then \(\operatorname{dist}(b,b_i)\ge2\rho'\) for
every \(i\). Let \(g\) be an isomorphism of the \(\rho'\)-neighbourhoods of
\(a\) and \(b\) with \(g(a)=b\), and let \(h'\) be the union of \(g\) and
the restriction of \(h\) to \(N(\bar a,\rho')\).

The sets \(D_1=N(\bar a,\rho')\) and \(D_2=N(a,\rho')\) are disjoint, and
no edge joins them: if \(x\in N(a_i,\rho')\) and \(y\in D_2\) were equal or
adjacent, then \(\operatorname{dist}(a,a_i)\le(\rho'-1)+1+(\rho'-1)=
2\rho'-1\). For the same reason, using \(\operatorname{dist}(b,b_i)\ge
2\rho'\), the images \(h(D_1)=N(\bar b,\rho')\) and \(g(D_2)=N(b,\rho')\)
are disjoint and joined by no edge. So \(h'\) is a bijection from
\(N(\bar aa,\rho')\) onto \(N(\bar bb,\rho')\); it preserves colours; it
preserves and reflects edges inside \(D_1\) (as \(h\) does) and inside
\(D_2\) (as \(g\) does); and there are no edges across on either side.
Thus \(h'\) is a local isomorphism with \(h'(a_i)=b_i\) and \(h'(a)=b\),
which is \((I_{j+1})\).

Playing in this way, Duplicator reaches \((I_r)\) and wins.
\end{proof}
\end{annexitem}

\begin{annexitem}{annex:p03-w03-p03-058-pinch}{thm:p03-w03-p03-058-pinch}{Pinching a coloured cycle}
\begin{proof}
Indices are taken modulo \(n\). Write \(c(i)\) for the colour of \(i\) and
\(\delta=q-p\). Equality of the windows of \(p\) and \(q\) says
\(c(p+e)=c(q+e)\) for \(|e|\le d\).

\emph{Item 1.} Let \(f(i)=i+1\), and let \(f_B(i)=i+1\) for
\(i\notin\{p,q\}\), \(f_B(p)=q+1\), \(f_B(q)=p+1\). The edges of
\(\mathbf B\) are exactly the sets \(\{i,f_B(i)\}\). The map \(f_B\) is a
permutation: it takes the values \(i+1\) for \(i\ne p,q\), which are all
vertices except \(p+1\) and \(q+1\), and these two are \(f_B(q)\) and
\(f_B(p)\). Its orbit through \(p+1\) is \(p+1,p+2,\ldots,q\), of length
\(\delta\), and its orbit through \(q+1\) is \(q+1,\ldots,n-1,0,\ldots,p\),
of length \(n-\delta\); both lengths are at least \(2d+1\ge3\). In an orbit
of length at least \(3\), each vertex \(i\) is adjacent exactly to
\(f_B(i)\) and \(f_B^{-1}(i)\), which are distinct and different from
\(i\). So \(E^{\mathbf B}\) is symmetric and irreflexive and each orbit
carries a cycle; no edge joins the two orbits, so \(\mathbf B\) is not
connected.

\emph{Neighbourhoods in long cycles.} Let \(F\) be a graph whose edges are
the sets \(\{i,g(i)\}\) for a permutation \(g\) all of whose orbits have
length \(L\ge2d+1\). For a vertex \(x\), the vertices \(g^j(x)\) with
\(|j|\le d-1\) are distinct (as \(2d-1<L\)), and \(g^j(x)\) is at distance
\(\min(|j|,L-|j|)=|j|\) from \(x\) (as \(|j|\le d-1<L/2\)); every other
vertex is at distance at least \(d\) or in another orbit. So
\(N(x,d)=\{g^j(x):|j|\le d-1\}\). Two of these, \(g^j(x)\) and
\(g^{j'}(x)\), are adjacent exactly when \(j'-j\equiv\pm1\pmod L\), and
since \(|j'-j|\le2d-2<L-1\) this means \(|j'-j|=1\). So the induced subgraph
on \(N(x,d)\) is the path \(g^{-(d-1)}(x),\ldots,g^{d-1}(x)\). This applies
to \(C_n\) with \(g=f\), \(L=n\), and to \(\mathbf B\) with \(g=f_B\).

\emph{Colours along the two paths.} We show \(c(f_B^j(x))=c(f^j(x))\) for
\(|j|\le d-1\). For \(j\ge0\), consider the walk
\(x,f_B(x),\ldots,f_B^j(x)\). If none of its first \(j\) points is \(p\) or
\(q\), then \(f_B^j(x)=x+j=f^j(x)\). Otherwise let \(m<j\) be the first index
with \(f_B^m(x)\in\{p,q\}\), so \(f_B^m(x)=x+m\). If \(x+m=p\), the next
point is \(q+1\), and the walk continues \(q+2,q+3,\ldots\) without meeting
\(p\) or \(q\) again within the remaining \(j-m-1<d\) steps, since it stays
in the orbit of \(q+1\), which returns to \(p\) only after
\(n-\delta-1\ge2d\) steps. So \(f_B^j(x)=q+(j-m)\) while
\(f^j(x)=p+(j-m)\), and \(1\le j-m\le d-1\) gives equal colours. If
\(x+m=q\), symmetrically \(f_B^j(x)=p+(j-m)\) and \(f^j(x)=q+(j-m)\),
the walk staying in the orbit of \(p+1\), which returns to \(q\) only after
\(\delta-1\ge2d\) steps; again the colours agree. For \(j<0\) use
\(f_B^{-1}\), which agrees with \(i\mapsto i-1\) except that
\(f_B^{-1}(q+1)=p\) and \(f_B^{-1}(p+1)=q\). If the backward walk from
\(x\) first reaches \(q+1\) after \(m\) steps, then
\(f_B^{-(m+1+e)}(x)=p-e\) and \(f^{-(m+1+e)}(x)=q-e\) for
\(0\le e\le d-2\); if it first reaches \(p+1\), the roles of \(p\) and \(q\)
are exchanged. In both cases the window equality gives equal colours.

\emph{Item 2.} For each vertex \(x\), the map \(f^j(x)\mapsto f_B^j(x)\)
(\(|j|\le d-1\)) is a bijection from \(N(x,d)\) in \(C_n\) onto \(N(x,d)\)
in \(\mathbf B\) that sends \(x\) to \(x\), preserves colours, and maps the
path onto the path. It is an isomorphism of the coloured
\(d\)-neighbourhoods, so \(x\) has the same \(d\)-type in both. The identity
of the common vertex set therefore preserves \(d\)-types, so for each
\(d\)-type the two coloured graphs have the same number of elements of
that type.
\end{proof}
\end{annexitem}

\begin{annexitem}{annex:p03-w03-p03-058-connectivity}{thm:p03-w03-p03-058-connectivity}{Connectivity is not monadic ESO-definable}
\begin{proof}
Fix \(k,r\ge1\) and let \(d=3^{r-1}\), \(M=2^{k(2d+1)}\),
\(n=(M+1)(2d+1)\). Let \(\mathcal K\) be any of the three combinations of a
class and a query in the statement. Duplicator plays as follows.

\emph{Choice of \(\mathbf A\).} She chooses the undirected cycle
\(\mathbf A=C_n\), whose edge relation is symmetric and irreflexive. Any two
vertices are joined by the path through consecutive vertices, which is a
directed path in either direction because \(E\) is symmetric. So
\(\mathbf A\) is connected and strongly connected, and the query takes the
value \(1\).

\emph{Choice of \(\mathbf B\).} Let Spoiler choose unary
\(S_1,\ldots,S_k\) on \(\{0,\ldots,n-1\}\). A window is a sequence of
\(2d+1\) colours from \(\{0,1\}^k\), so there are \(M\) windows. By the
pigeonhole principle two of the \(M+1\) vertices \(m(2d+1)\),
\(0\le m\le M\), have the same window; call them \(p<q\). Then
\(q-p\ge2d+1\), and \(q\le M(2d+1)=n-(2d+1)\) gives
\(n-(q-p)\ge n-q\ge2d+1\). Duplicator chooses the graph \(\mathbf B\) of
\cref{thm:p03-w03-p03-058-pinch} and \(S'_i=S_i\). By that lemma
\(\mathbf B\) has symmetric irreflexive edge relation and is not connected;
being symmetric, it is not strongly connected either. So \(\mathbf B\)
belongs to the class and the query takes the value \(0\).

\emph{The first-order stage.} \((\mathbf A,S_1,\ldots,S_k)\) and
\((\mathbf B,S_1,\ldots,S_k)\) are finite \(k\)-coloured graphs, and by
\cref{thm:p03-w03-p03-058-pinch} they are \(d\)-equivalent. Since
\(d=3^{r-1}\), \cref{thm:p03-w03-p03-058-hanf} gives Duplicator a winning
strategy in the \(r\)-move game on them.

So Duplicator wins the \((k,r)\)-Ajtai--Fagin game, for every \(k\) and
\(r\). By \cref{thm:p03-w03-p03-058-af-method}(1) the query is not monadic
ESO-definable on the class.
\end{proof}
\end{annexitem}

\begin{annexitem}{annex:p03-w03-p03-058-complement}{thm:p03-w03-p03-058-complement}{Monadic ESO is not closed under complement}
\begin{proof}
Let \(\mathbf G=(V,E)\) be a finite graph.

\emph{\(\delta\) and strong connectivity.} If \(\mathbf G\) is not
strongly connected, there are \(u,v\) with no directed path from \(u\) to
\(v\). Let \(S\) be the set of vertices reachable from \(u\) by a directed
path (including \(u\)). Then \(u\in S\) and \(v\notin S\), and if
\(z\in S\) and \(E(z,w)\) then \(w\in S\); so \(S\) witnesses \(\delta\).
Conversely, let \(S\) witness \(\delta\), with \(x\in S\) and \(y\notin S\).
A directed path from \(x\) to \(y\) would contain a first edge \((z,w)\)
with \(z\in S\) and \(w\notin S\), which the matrix forbids. So \(\mathbf G\)
is not strongly connected. Thus \(\delta\) defines the failure of strong
connectivity on all finite graphs.

\emph{\(\delta'\) and connectivity.} If \(\mathbf G\) is not connected,
let \(S\) be the set of vertices joined to a fixed \(u\) by a path in the
symmetric closure of \(E\); some \(v\) lies outside \(S\), and an edge in
either direction between \(z\in S\) and \(w\notin S\) would join \(w\) to
\(u\). So \(S\) witnesses \(\delta'\). Conversely, if \(S\) witnesses
\(\delta'\) with \(x\in S\), \(y\notin S\), a path from \(x\) to \(y\) in the
symmetric closure has a first step between a vertex of \(S\) and one
outside, which is an edge in one direction or the other, and the matrix of
\(\delta'\) forbids both. So \(\delta'\) defines disconnectivity on all
finite graphs.

\emph{Symmetric graphs.} If \(E\) is symmetric, \(\neg E(z,w)\) and
\(\neg E(z,w)\wedge\neg E(w,z)\) are equivalent, so \(\delta\) and
\(\delta'\) agree and \(\delta\) defines disconnectivity on finite graphs
with symmetric irreflexive edge relation.

In each case the complementary query, connectivity or strong connectivity on
the class concerned, is not monadic ESO-definable by
\cref{thm:p03-w03-p03-058-connectivity}. Hence a monadic ESO-definable query
has a complement that is not monadic ESO-definable.
\end{proof}
\end{annexitem}

\begin{annexitem}{annex:p03-w03-p03-058-reduct}{thm:p03-w03-p03-058-reduct}{Adding relations to the vocabulary}
\begin{proof}
Let \(\Psi=(\exists P_1)\cdots(\exists P_k)\psi\) define \(Q\) on
\(\mathcal C\), with \(\psi\) first-order over
\(\sigma\cup\{P_1,\ldots,P_k\}\). The relation variables are chosen distinct
from the symbols of \(\sigma'\) (rename them if necessary), so \(\Psi\) is
also an ESO-sentence over \(\sigma'\), monadic when it was monadic over
\(\sigma\).

Let \(\mathbf G\in\mathcal C'\) with universe \(V\) and put
\(\mathbf G_0=\mathbf G{\restriction}\sigma\in\mathcal C\). For relations
\(S_1,\ldots,S_k\) on \(V\) of the right arities, the structures
\((\mathbf G,S_1,\ldots,S_k)\) and \((\mathbf G_0,S_1,\ldots,S_k)\) have the
same universe and interpret every symbol of
\(\sigma\cup\{P_1,\ldots,P_k\}\) alike. By induction on first-order
formulas over that vocabulary, an assignment satisfies a formula in one
structure exactly when it satisfies it in the other: atomic formulas are
evaluated by the common interpretations, the connectives commute with the
induction hypothesis, and the quantifiers range over the common universe
\(V\). Hence \((\mathbf G,S_1,\ldots,S_k)\entails\psi\) iff
\((\mathbf G_0,S_1,\ldots,S_k)\entails\psi\), and since the witnesses range
over the same relations on \(V\) in both cases,
\(\mathbf G\entails\Psi\) iff \(\mathbf G_0\entails\Psi\). Therefore
\[
 Q'(\mathbf G)=1\iff Q(\mathbf G_0)=1\iff\mathbf G_0\entails\Psi\iff
 \mathbf G\entails\Psi,
\]
and \(\Psi\) defines \(Q'\) on \(\mathcal C'\).

For the last claim take \(\sigma=\{E\}\), \(\mathcal C\) the class of all
finite graphs, and \(\sigma'=\{E,\mathrm{Suc}\}\) with \(\mathcal C'\) the
finite structures with successor, or \(\sigma'=\{E,\le\}\) with
\(\mathcal C'\) the finite ordered graphs. The \(\{E\}\)-reduct
\((V,E)\) of a member of \(\mathcal C'\) is a finite graph, and the
Connectivity query on \(\mathcal C'\) is the query \(Q'\) obtained from the
Connectivity query \(Q\) on finite graphs, in the same reading, since its
value depends on \((V,E)\) alone
(\cref{def:p03-w03-p03-058-builtin}). If \(Q\) were monadic ESO-definable on
\(\mathcal C\), the first part would make \(Q'\) monadic ESO-definable on
\(\mathcal C'\), contrary to \cref{thm:p03-w03-p03-058-successor} or
\cref{thm:p03-w03-p03-058-ordered} respectively.
\end{proof}
\end{annexitem}

\begin{annexitem}{annex:p03-w03-p03-058-binary-conn}{thm:p03-w03-p03-058-binary-conn}{Connectivity in binary ESO}
\begin{proof}
Write \(\lambda(L)\) as the conjunction of \((\forall x)\neg L(x,x)\),
\((\forall x)(\forall y)(\forall z)(L(x,y)\wedge L(y,z)\to L(x,z))\) and
\((\forall x)(\forall y)(x=y\vee L(x,y)\vee L(y,x))\). Let \(\mathbf G\) be a
finite \(\sigma\)-structure with universe \(V\ne\emptyset\). For a strict
linear order \(L\) on \(V\), the condition \(\neg(\exists w)L(w,x)\) says
that \(x\) is the \(L\)-least element; since \(V\) is finite and nonempty
there is exactly one such element, written \(m_L\). A \emph{path} below
means a sequence \(v_0,\ldots,v_\ell\) (\(\ell\ge0\)) with
\(E(v_{i-1},v_i)\vee E(v_i,v_{i-1})\) for each \(i\) (undirected), or with
\(E(v_{i-1},v_i)\) for each \(i\) (directed).

\emph{\(\gamma\) implies connectivity.} Let \(L\) witness \(\gamma\) and
\(m=m_L\). We show by induction along \(L\), which is well founded because
\(V\) is finite, that every \(y\in V\) is joined to \(m\) by an undirected
path. For \(y=m\) take the path \(m\). For \(y\ne m\) we have \(L(m,y)\), so
the second conjunct of \(\gamma\) gives \(z\) with \(L(z,y)\) and
\(E(z,y)\vee E(y,z)\). By the induction hypothesis \(z\) is joined to
\(m\); appending \(y\) gives an undirected path from \(m\) to \(y\). Any two
vertices are then joined through \(m\), so \((V,E)\) is connected.

\emph{Connectivity implies \(\gamma\).} Fix \(m\in V\) and let
\(\mathrm{dist}(y)\) be the least length of an undirected path from \(m\)
to \(y\), finite for every \(y\) by connectivity. List \(V\) as
\(v_0,\ldots,v_{n-1}\) with \(\mathrm{dist}\) non-decreasing and
\(v_0=m\) (only \(m\) has distance \(0\)), and let \(L(v_i,v_j)\) iff
\(i<j\); this is a strict linear order with least element \(m\). If
\(y\ne m\) and \(\mathrm{dist}(y)=\ell\ge1\), the vertex \(z\) preceding
\(y\) on a shortest path has \(\mathrm{dist}(z)=\ell-1\), and it is
adjacent to \(y\) in the symmetric closure. As
\(\mathrm{dist}(z)<\mathrm{dist}(y)\), \(z\) precedes \(y\) in the list, so
\(L(z,y)\). Every \(y\) with \((\exists w)L(w,y)\) is such a \(y\ne m\), so
\(L\) witnesses \(\gamma\).

\emph{\(\gamma_s\) implies strong connectivity.} Let \(L_1,L_2\) witness
\(\gamma_s\). The third conjunct makes \(m_{L_1}=m_{L_2}\); call it \(m\).
By induction along \(L_1\) as above, using the conjunct for \(L_1\), every
\(y\) is reached from \(m\) by a directed path: for \(y\ne m\) there is
\(z\) with \(L_1(z,y)\) and \(E(z,y)\), and a directed path from \(m\) to
\(z\) followed by the edge \((z,y)\) reaches \(y\). By induction along
\(L_2\), using the conjunct for \(L_2\), every \(y\) reaches \(m\) by a
directed path: for \(y\ne m\) there is \(z\) with \(L_2(z,y)\) and
\(E(y,z)\), and the edge \((y,z)\) followed by a directed path from \(z\)
to \(m\) reaches \(m\). For any \(u,v\), a directed path from \(u\) to
\(m\) followed by one from \(m\) to \(v\) is a directed path from \(u\) to
\(v\). So \((V,E)\) is strongly connected.

\emph{Strong connectivity implies \(\gamma_s\).} Fix \(m\in V\). Let
\(d_1(y)\) be the least length of a directed path from \(m\) to \(y\) and
\(d_2(y)\) that of a directed path from \(y\) to \(m\); both are finite.
Build \(L_1\) from \(d_1\) and \(L_2\) from \(d_2\) as in the second step,
each with least element \(m\), the only vertex at distance \(0\). The third
conjunct holds because both orders have least element \(m\). For
\(y\ne m\), the vertex \(z\) preceding \(y\) on a shortest directed path
from \(m\) has \(E(z,y)\) and \(d_1(z)=d_1(y)-1\), so \(L_1(z,y)\); the
vertex \(z\) following \(y\) on a shortest directed path from \(y\) to
\(m\) has \(E(y,z)\) and \(d_2(z)=d_2(y)-1\), so \(L_2(z,y)\). Thus
\(L_1,L_2\) witness \(\gamma_s\).

Only the binary relation variables \(L\), \(L_1\), \(L_2\) are quantified,
and no symbol of \(\sigma\) other than \(E\) occurs, so the equivalences
hold on every finite \(\sigma\)-structure, in particular on finite graphs,
on finite structures with successor and on finite ordered graphs.
\end{proof}
\end{annexitem}
